\section{Threats to Validity}
\textbf{Construct validity.} Typed topology and explicit rules do not measure responsibility, rationale, runtime quality, or developer usefulness. Evidence agreement requires one matching location, not complete or well-ranked evidence. An absent relationship has only a deterministic component anchor. D2 signals can share an edge, and its negatives do not represent all unsupported syntax. Parsed-file counts measure analysis scope, not CPU work.

\textbf{Internal validity.} The authors developed the analyzer and both datasets. D1 topology intent preceded integration, but patch implementations and exact evidence lines co-evolved with the analyzer. D2 targeted known v0.1 errors and guided v0.2. Replaying D1 in P3 adds incremental checks, not independent observations. Hashes, raw predictions, mutation checks, and duplicate execution expose inconsistency but do not establish independent ground truth.

\textbf{External and conclusion validity.} One small synthetic system, 20 patches, and 40 curated signals cannot estimate general TypeScript accuracy. The five changed nodes, 12 changed edges, seven violation rule sets, and 11 evidence cases are finite regression items. Wrappers, dynamic endpoints, dependency injection, generated code, unconventional repositories, other languages, and runtime-only relationships remain outside evaluated coverage. Latency is descriptive from one environment; there is no workload-population inference, scalability study, or significance test. No industrial deployment or developer study was executed. The private repository plan did not enforce non-bypassable merge protection.

\textbf{Comparative validity.} The executed external inventory has no common labeled units and supplies no accuracy baseline. Unresolved imports reflect uninstalled subject dependencies rather than labeled analyzer errors. Distinct observation semantics and newer pinned packages prevent its import inventory from extending the historical service-graph accuracy or timing claims. Unsupported service edges are excluded rather than counted as comparator misses. Independently authored data and capability-matched ground truth remain necessary.

\textbf{Literature positioning.} The narrative review carries source-selection and interpretation bias and may omit approaches. Coverage completeness and absence-of-evidence claims are outside its design, not deferred SLR deliverables. Related studies can share projects or implementations, so citation breadth does not imply independent replication.

\textbf{Reliability.} Pinned artifacts, normalized ordering, isolated patches, duplicate execution, and clean Linux/Windows functional checks support reproducibility. The frozen v0.1 result is provenance-validated rather than continuously regenerated; Windows timing is not a cross-platform performance claim. The external run's invalid first attempt remains auditable and excluded, with the corrected patch-isolation procedure frozen before rerunning.
